\begin{abstract}
Natural language is emerging as a primary feedback channel for improving language agents, capable of conveying intent, preferences, and causal structure in forms interpretable by both humans and modern language models. We call this paradigm Verbal Reinforcement Learning (VRL) and offer the first unified account of it. We organize the field around a single axis, \textit{when} verbal feedback takes effect in an agent's lifecycle and \textit{what} it modifies, yielding three pillars: (1) \textbf{Language as Grounding Signal}, where language defines the task itself by specifying goals, states, and reward structures; (2) \textbf{Language as Deliberative Feedback}, where natural language guides reasoning at test time without the need to update model parameters; (3) \textbf{Language as Learning Signal}, where language-based feedback shapes model parameters through training. Within each pillar, we synthesize representative work, distinguish key subcategories of approaches, and outline the distinct role language plays in shaping agent behavior. Together, this taxonomy shows how verbal reinforcement is reshaping agent development, while also defining the challenges and opportunities for building more capable and aligned agents.
\end{abstract}

% Through our analysis, we identify common patterns and introduce a taxonomy based on the roles language plays:
% This survey synthesizes the rapidly expanding literature on
% Verbal Reinforcement Learning (VRL): a paradigm where natural language feedback drives the improvement of language agents. Such feedback is uniquely expressive, conveying intent, preferences, and causal structure in a form interpretable by both humans and modern language models.  Its rise reflects a broader pattern in machine learning: just as self-supervised learning unlocked the potential of unlabeled data, verbal feedback is now emerging as a key driver of progress in language agents, and this survey synthesizes the rapidly expanding literature mapping that shift. 